% Part: many-valued-logic
% Chapter: syntax-and-semantics
% Section: introduction

\documentclass[../../../include/open-logic-section]{subfiles}

\begin{document}

\olfileid{mvl}{syn}{int}

\olsection{ভূমিকা}

ধ্রুপদি যুক্তিবিদ্যায় আমরা এমন !!{formula}s নিয়ে কাজ করি, যেগুলি
!!{propositional variable}s থেকে বচনসংযোজক $\lnot$, $\land$, $\lor$,
$\lif$ ও~$\liff$ ব্যবহার করে গঠিত। ধ্রুপদি যুক্তিবিদ্যার অর্থতত্ত্ব
সংজ্ঞায়িত করার সময় আমরা $\True$ ও~$\False$—এই দুই সত্যমান ব্যবহার করি।
!!{propositional variable}s-গুলিকে আমরা !!a{valuation}~$\pAssign{v}$-এ
ব্যাখ্যা করি; সেটি ওই !!{propositional variable}s-কে $\True$, $\False$
সত্যমান আরোপ করে। তখন যেকোনো !!{valuation} প্রতিটি !!{formula}~$!A$-এর
একটি সত্যমান $\pValue{v}(!A)$ নির্ধারণ করে, এবং !!^a{formula} কোনো
!!a{valuation}~$\pAssign{v}$-এ পরিতৃপ্ত হয়, $\pSat{v}{!A}$, তখনই যখন
$\pValue{v}(!A) = \True$।

বহুমানী যুক্তিবিদ্যা ধ্রুপদি দ্বিমানী যুক্তিবিদ্যার এমন সাধারণীকরণ, যেখানে
$\True$ ও~$\False$ ছাড়াও আরও সত্যমান অনুমোদিত। তাই বহুমানী যুক্তিবিদ্যায়
!!a{valuation}~$\pAssign{v}$ এমন একটি অপেক্ষক, যা প্রতিটি
!!{propositional variable}~$p$-কে সম্ভাব্য সত্যমানগুলির কোনো একটি আরোপ
করে। অনুমোদিত সত্যমানের সেটকে আমরা সাধারণত~$V$ বলব। ধ্রুপদি যুক্তিবিদ্যাও
একটি বহুমানী যুক্তিবিদ্যা, যেখানে $V = \{\True, \False\}$; আর সংযোজকগুলির
পরিচিত বৈশিষ্ট্যসূচক সত্যসারণি ব্যবহার করে সত্যমান~$\pValue{v}(!A)$
গণনা করা হয়।

অতিরিক্ত সত্যমান যোগ করলে “এবং”, “অথবা”, “নয়” ও “যদি---তবে” হিসেবে
পড়া সংযোজকগুলির জন্য $\pValue{v}(!A)$ গণনার একাধিক স্বাভাবিক উপায়
পাওয়া যায়। তাই কোনো বহুমানী যুক্তিবিদ্যা শুধু তার সত্যমানের সেট দিয়েই
নির্ধারিত হয় না; প্রতিটি সংযোজকের জন্য আমরা যে \emph{সত্যমান-অপেক্ষক}
বেছে নিই, সেগুলিও সেটিকে নির্ধারণ করে। কিন্তু কোনো বহুমানী
যুক্তিবিদ্যা~$\Log L$-এর জন্য এগুলি একবার বেছে নিলে, ধ্রুপদি যুক্তিবিদ্যার
মতোই সত্যমান-আরোপ থেকে সত্যমান $\pValue{v}(!A)[\Log L]$ এককভাবে
নির্ধারিত হয়। ধ্রুপদি যুক্তিবিদ্যার মতো বহুমানী যুক্তিবিদ্যাও
\emph{সত্যমান-অপেক্ষকধর্মী}।

এই অর্থগত গঠন-উপাদানগুলি হাতে থাকলে সর্বতঃসত্য, অনুসিদ্ধান্ত ও
পরিতৃপ্তিযোগ্যতার অনুরূপ ধারণাগুলি সংজ্ঞায়িত করা যায়। ধ্রুপদি
যুক্তিবিদ্যায় কোনো !!a{formula} সর্বতঃসত্য, যদি প্রতিটি~$\pAssign{v}$-এর
জন্য তার সত্যমান $\pValue{v}(!A) = \True$ হয়। বহুমানী যুক্তিবিদ্যায়
এটিকেও কিছুটা সাধারণীকরণ করতে হয়। প্রথমত, সত্যমানের সেট~$V$-তে
$\True$ থাকতেই হবে এমন কোনো শর্ত নেই। যেমন, কোনো কোনো বহুমানী
যুক্তিবিদ্যায় $0$ ও~$1$-এর মধ্যবর্তী সব মূলদ সংখ্যাকে সত্যমানের সেট
হিসেবে নেওয়া হয়। সে ক্ষেত্রে $1$ সাধারণত~$\True$-এর ভূমিকা পালন করে।
অন্য কোনো যুক্তিবিদ্যায় একটির বদলে একাধিক সত্যমান সেই ভূমিকা পালন করে।
তাই প্রতিটি বহুমানী যুক্তিবিদ্যার \emph{মনোনীত মানসমূহ}-এর একটি
সেট~$V^+$ থাকা চাই। এরপর বলা যায়, কোনো !!a{formula}
!!a{valuation}~$\pAssign{v}$-এ পরিতৃপ্ত, $\pSat{v}{!A}[\Log L]$, তখনই
যখন $\pValue{v}(!A)[\Log L] \in V^+$। কোনো !!^a{formula}~$!A$ ওই
যুক্তিবিদ্যার সর্বতঃসত্য, $\Entails[\Log L] !A$, তখনই যখন প্রতিটি
$\pAssign{v}$-এর জন্য $\pValue{v}(!A) \in V^+$। এবং সব শেষে, আমরা বলি
!!{formula}s-এর কোনো সেট থেকে~$!A$ অনুসৃত হয়,
$\Gamma \Entails[\Log L] !A$, যদি~$\Gamma$-র সব !!{formula}s-কে
পরিতৃপ্ত করে এমন প্রত্যেক !!{valuation}~$!A$-কেও পরিতৃপ্ত করে।

\end{document}
